% Build: latexmk -cd -pdf doc/project2/project2_solution.tex
\documentclass[a4paper,10pt]{article}
\usepackage[a4paper,left=15mm,right=15mm,top=15mm,bottom=20mm]{geometry}
\usepackage[utf8]{inputenc}
\usepackage{kotex}
\usepackage{amsmath,amssymb,mathtools}
\usepackage{enumitem,booktabs}
\usepackage[hidelinks]{hyperref}

\newcommand{\ps}[1]{\mathcal P(#1)}
\newcommand{\addr}{\operatorname{addr}}
\newcommand{\Msharp}{M^\sharp}
\newcommand{\Rsharp}{R^\sharp}
\newcommand{\Isharp}{I^\sharp}
\newcommand{\Csharp}{C^\sharp}
\newcommand{\Ssharp}{S^\sharp}
\newcommand{\remN}[1]{\overline{#1}}

\title{프로그램분석 26F Project 2\\[2pt]
  \large Abstract domain 제안}
\author{2024-21504 남연우(Yeonwoo Nam)}
\date{}

\begin{document}
\maketitle

\paragraph{범위.}
현재 문서는 워터첵 분석기에 사용할 요약 의미공간을 제안하는 초안이다.
구체적인 요약 의미연산자 정의와 그 안전성 증명은 이후에 작성한다.
분석 방식은 Project~1의 transition style을 사용하여,
종료 시점뿐 아니라 모든 도달 가능한 프로그램 지점의 메모리를 추적한다.

\paragraph{기본 설정.}
$N=260401$로 두고, 분석할 프로그램에 등장하는 변수들의 유한집합을 $V$라 하자.
Project~1과 같이 메모리는 $m:\mathbb Z\to\mathbb Z$이고,
변수 $x$의 값은 $m(\addr(x))$이다.
변수 주소 함수 $\addr$는 단사이며, 추적하지 않는 주소의 메모리도 존재한다.
성질 P1과 P2의 대상은 $V$에 속한 변수들이고, P2에서는 서로 다른 두 변수를 고려한다.
초기 메모리집합 $M_0$는 주어진 매개변수로 유지한다.

\paragraph{제안하는 요약 메모리.}
P1을 위해 변수별 나머지를 추적하고,
P2를 위해 변수 및 두 변수의 합과 차에 대한 interval 및 congruence 정보를 추적한다.
구체적으로 다음 곱공간을 사용한다:
\[
 \boxed{
 \Msharp=(V\to\Rsharp)\times(T\to\Ssharp),
 \qquad \Ssharp=\Isharp\times\Csharp.
 }
\]
여기서 $T$는 추적할 식(template)의 유한집합이다:
\[
 T=\{x\mid x\in V\}
   \cup\{x+y,\ x-y\mid x,y\in V,\ x\prec y\}.
\]
$\prec$는 같은 변수 쌍을 한 번만 저장하기 위한 임의의 고정 전순서이다.
요약 메모리 원소를 $a=(\rho,\kappa)$로 쓴다.

\paragraph{1. P1: 나머지 집합.}
$\mathbb Z_N=\{0,\ldots,N-1\}$로 두고,
\[
 \Rsharp=\ps{\mathbb Z_N},\qquad
 R\sqsubseteq_R S\iff R\subseteq S
\]
로 정의한다. $\remN{z}=z\bmod N\in\mathbb Z_N$로 표기하면,
$\rho(x)=R$의 의미는
\[
 \remN{m(\addr(x))}\in R
\]
이다. 최하원소는 $\emptyset$, 최상원소는 $\mathbb Z_N$이고,
합류점에서는 가능한 나머지들의 합집합을 취한다.
예를 들어 $\rho(x)=\{0,1\}$이면 정확한 정숫값은 모르더라도
$x$가 P1을 만족하지 않는다는 것은 알 수 있다.

산술은 나머지에 대해 계산할 수 있다. 예를 들어 덧셈은
\[
 R+^\sharp S=\{\remN{r+s}\mid r\in R,\ s\in S\}
\]
로 근사한다. 상수 $n$은 $\{\remN{n}\}$,
\texttt{readInt}는 $\mathbb Z_N$으로 근사한다.

\paragraph{2. P2: template별 interval과 congruence.}
초기 제안의 flat constant domain은 서로 다른 두 상수의 join이 곧 $\top$이다.
따라서 합이 변하지만 $123$은 될 수 없는 경우를 자주 놓친다.
P2의 부재를 증명하려면 합이 하나의 상수라는 사실보다
\emph{합이 가질 수 없는 값}에 대한 정보를 더 오래 유지해야 한다.
이를 위해 각 template의 값을 다음 두 성분으로 요약한다.

\begin{itemize}[leftmargin=*,itemsep=3pt]
 \item $\Isharp$: 정수 interval $[\ell,u]\cap\mathbb Z$와 공집합.
 하한에는 $-\infty$, 상한에는 $+\infty$를 허용한다.
 \item $\Csharp$: 정수 congruence $d\mathbb Z+r$와 공집합.
 여기서 $d\ge0$이며, $d=0$이면 singleton $\{r\}$,
 $d=1$이면 모든 정수 $\mathbb Z$를 나타낸다.
\end{itemize}

$\kappa(t)=(I_t,C_t)$의 의미는 다음과 같다:
\[
 t(m)\in\gamma_S(I_t,C_t)
       :=\gamma_I(I_t)\cap\gamma_C(C_t).
\]
즉, 가능한 값이 interval 안에 있으면서 congruence도 만족한다.
최상원소는 $([-\infty,+\infty],\mathbb Z)$이고,
상수 $c$는 $([c,c],0\mathbb Z+c)$로 표현할 수 있다.
이와 같이 interval과 congruence를 결합한 공간은
strided interval로 볼 수 있다. 두 성분 사이의 정보 교환(reduction)은
추후 연산자 설계에서 추가할 수 있다.\footnote{Interval과 congruence의
곱 및 reduction에 대해서는 Antoine Min\'e,
\emph{Tutorial on Static Inference of Numeric Invariants by Abstract Interpretation},
\S6.2를 참고한다:
\url{https://perso.lip6.fr/Antoine.Mine/publi/article-mine-FTiPL17.pdf}.}

차이는 다음 예들에서 드러난다. 아래는 한 지점의 특정 template에 대한 예이다.
\begin{center}
\begin{tabular}{lll}
\toprule
가능한 합 & Flat constants & Interval $\cap$ congruence \\
\midrule
$\{100,110\}$ & $\top$ & $[100,110]\cap10\mathbb Z$ \\
$\{122,124\}$ & $\top$ & $[122,124]\cap2\mathbb Z$ \\
$2\mathbb Z$ & $\top$ & $\mathbb Z\cap2\mathbb Z$ \\
\bottomrule
\end{tabular}
\end{center}
세 경우 모두 새로운 표현에서는 $123$을 배제할 수 있다.
특히 $122$와 $124$의 join에서 interval만 사용하면 $123$이 포함되지만,
congruence 성분은 짝수라는 사실을 유지한다.

\paragraph{Congruence가 P2에 주는 정보.}
$C_{x+y}=d\mathbb Z+r$이고 $d>0$이면,
\[
 x+y=123\quad\Longrightarrow\quad d\mid(123-r).
\]
따라서 $d\nmid(123-r)$이면 이 변수 쌍은 P2를 만족할 수 없다.
$d=0$인 singleton의 경우에는 $r\ne123$인지 검사한다.
반대로 congruence에 $123$이 포함된다는 사실만으로
P2가 실제로 발생한다고 결론낼 수는 없다.

피연산자를 독립적으로 요약하여 계산하는 표준적인 식 평가에서도
congruence는 유용하다. 예를 들어 다음 명령들을 생각하자:
\begin{quote}\ttfamily
 x := 2 * readInt;\\
 y := 2 * readInt
\end{quote}
각 변수에 대해 $C_x=C_y=2\mathbb Z$를 얻으므로,
독립적인 덧셈 평가만으로도
\[
 C_x+^\sharp C_y=2\mathbb Z+2\mathbb Z=2\mathbb Z
\]
를 얻는다. 두 interval이 모두 무한 범위여도 합은 짝수이므로 $123$을 배제한다.
이 예는 마지막 지점의 한 변수 쌍에 관한 것이며,
프로그램 전체의 P1 또는 P2 부재를 뜻하지는 않는다.

\paragraph{도메인의 표현력과 연산자의 정확도.}
Template을 추가하는 것만으로 변수 사이의 상관관계를 자동으로 추론하지는 않는다.
예를 들어 \texttt{x := readInt; y := x} 직후에
각 변수의 congruence는 $C_x=C_y=\mathbb Z$이다.
독립적인 덧셈 평가는
\[
 C_x+^\sharp C_y=\mathbb Z+\mathbb Z=\mathbb Z
\]
를 내놓으며, 이는 안전한 근사이다.
실제로는 $x+y=2x\in2\mathbb Z$이지만,
이 사실은 앞의 독립적인 계산만으로 얻어지지 않는다.
제안한 공간은 이를 $x+y$ template에 표현할 수 있다.
다만 그 정보를 얻으려면 \texttt{y := x}의 대입 연산자가
$(x+y)[y\leftarrow x]=x+x=2x$처럼 치환 및 affine 정규화를 수행하도록
별도로 정의하고 안전성을 확인해야 한다.
$x-y$ template도 복사 관계와 상대적인 차이를 표현하기 위해 포함한 것이며,
그 관계를 유지하는 연산자는 별도 설계 대상이다.

Interval 성분은 $\pm x\pm y\le c$ 형태의 제약을 표현하므로,
octagon 계열의 관계형 추상화로 볼 수 있다.\footnote{Antoine Min\'e,
\emph{The Octagon Abstract Domain}:
\url{https://mine.perso.lip6.fr/publi/article-mine-HOSC06.pdf}.}
여기에 template별 congruence를 추가하여 범위뿐 아니라
짝홀성, 일정 간격의 값들, 정확한 상수도 함께 표현한다.
Octagon의 closure나 관계형 대입 연산자가 이 공간 정의만으로 주어지는 것은 아니다.

\paragraph{곱공간의 해석과 순서.}
요약 메모리의 의미는 두 성분의 제약을 동시에 만족하는 메모리들이다:
\[
 \gamma_M(\rho,\kappa)=
 \left\{m\ \middle|\
 \begin{aligned}
  &\forall x\in V.\ \remN{m(\addr(x))}\in\rho(x),\\
  &\forall t\in T.\ t(m)\in\gamma_S(\kappa(t))
 \end{aligned}
 \right\}.
\]
여기서 $x(m)=m(\addr(x))$이고
$(x+y)(m)=m(\addr(x))+m(\addr(y))$이다.
추적하지 않는 메모리 셀에는 제약을 두지 않는다.
Interval과 congruence 각각은 나타내는 정수 집합의 포함관계로 순서를 정하고,
전체 곱공간에는 점별 곱순서를 사용한다.
Join은 나머지 집합의 합집합, interval의 convex hull,
두 congruence를 모두 포함하는 최소 congruence를 성분별로 계산한다.
Abstraction은 각 변수의 가능한 나머지를 모으고,
각 template의 가능한 값들을 포함하는 최소 interval과 최소 congruence로 요약한다.

\paragraph{워터마크 부재 판정.}
각 프로그램 지점의 요약 메모리에 대해 다음 충분조건을 검사한다:
\begin{itemize}[leftmargin=*,itemsep=3pt]
 \item P1 부재:
 $\forall x\in V.\ 123\notin\rho(x)$.
 \item P2 부재: 모든 서로 다른 변수 쌍 $x,y$에 대해,
 $123\notin\gamma_S(\kappa(x+y))$이거나,
 $123\notin\rho(x)+^\sharp\rho(y)$이다.
 여기서 $\kappa(x+y)$는 순서 $\prec$에 따라 정규화한 쌍의 항목을 뜻한다.
\end{itemize}
두 번째 항목의 나머지 검사는 $x+y=123$이면
$\remN{x+y}=123$이라는 사실을 사용한다.
$\kappa(x+y)=(I_{x+y},C_{x+y})$일 때 첫 번째 검사는
\[
 123\notin\gamma_I(I_{x+y})
 \quad\lor\quad
 123\notin\gamma_C(C_{x+y})
\]
와 같다. 즉, 범위 또는 congruence 중 어느 하나가 $123$을 배제하면 충분하다.
주어진 성질의 부재 조건이 모든 도달 가능한 지점에서 성립하면
그 성질이 실행 중 발생하지 않음을 결론낼 수 있다.
두 종류 모두를 배제하려면 P1과 P2를 각각 배제해야 한다.
조건이 성립하지 않는 경우에는 명세에 따라 결론을 내리지 않는다.

\paragraph{간접 메모리 접근.}
단일 변수 template은 포인터 해석에도 사용할 수 있다.
예를 들어 \texttt{i := \&x} 이후 다음 정보를 유지할 수 있다:
\[
 \gamma_S(\kappa(i))=\{\addr(x)\}.
\]
이 경우 \texttt{*i}를 $x$의 읽기로,
\texttt{*i := E}를 $x$의 대입으로 처리할 수 있다.
주소를 확정할 수 없는 간접 읽기는 알 수 없는 값으로,
간접 쓰기는 수정될 수 있는 추적 셀의 정보를 잊는 방식으로 보수적으로 처리한다.
특히 나머지만 같다는 이유로 두 주소를 같은 셀로 취급해서는 안 된다.

\paragraph{선택 이유와 한계.}
$T$는 유한하지만 interval 성분에는
$[0,0]\subsetneq[0,1]\subsetneq\cdots$ 같은 무한 상승 체인이 있다.
따라서 이전 flat constant 제안과 달리 전체 공간은 유한 높이가 아니며,
반복 계산의 종료를 위해 interval 성분에 widening 등을 도입해야 한다.
이때 interval이 넓어져도 congruence 성분은 별도로 유지하여
짝홀성과 같은 정보를 보존하는 방향으로 설계한다.
최종 결과가 초기 상태를 포함하고 모든 전이에 대해 닫힌 invariant임을 확인하면
Project~1의 transition-style 안전성 정리를 사용할 수 있다.

이 공간도 모든 비선형 관계나 경로별 disjunction을 보존하지는 않는다.
예를 들어 합이 $\{122,125,126\}$ 중 하나이면
interval과 congruence만으로는 빠진 값 $123$을 보존하지 못한다.
또한 관계가 없는 두 임의의 입력 변수의 합은 실제로 모든 정수가 가능하므로,
그 경우의 $\top$은 적절한 결과이다.
나머지 집합은 변수마다 최대 $N$비트가 필요하다.
더 작은 첫 구현이 필요하면 나머지 성분도
$\bot$, 단일 나머지, $\top$으로 이루어진 flat domain으로 바꿀 수 있다.
이 경우 서로 다른 두 나머지를 합류할 때의 정보 손실은 더 커진다.
초기 메모리를 완전히 임의로 허용하면 시작점부터 워터마크가 가능할 수 있으므로,
부재 판정의 유용성은 주어진 $M_0$에도 의존한다.

\paragraph{AI 활용 및 확인 범위.}
이 초안의 도메인 제안과 문서 작성에 AI를 사용하였다.
Project~2 명세 및 기존 Project~1의 메모리 모델과 대조하여 작성하였다.
현재 단계에서는 이 도메인의 요약 연산자나 Rocq 인스턴스를 구현하거나
그 안전성을 기계 검증하지 않았다.

\end{document}
